\documentclass[11pt,a4paper]{article}

\usepackage[utf8]{inputenc}
\usepackage[T1]{fontenc}
\usepackage{lmodern}
\usepackage{amsmath,amssymb,amsfonts,bm}
\usepackage{geometry}
\usepackage{hyperref}
\usepackage{authblk}
\usepackage{booktabs}
\usepackage{enumitem}

\geometry{margin=2.4cm}

\title{\textbf{Technical Note (v49.1)}\\[0.2cm]
\textbf{Numerical Evaluation of the Relaxation Scale}\\[0.15cm]
\large \(k_{\mathrm{rel}}=a_{\mathrm{fr}}H(a_{\mathrm{fr}})\)
from the Companion-XI background}

\author[1]{Michael Lehmann}
\affil[1]{\small Independent Researcher, Relaxation-Driven Cyclic Cosmology (RDCC)\\
\texttt{mi.lehmann@gmx.de}}

\date{September 2026}

\begin{document}
\maketitle

\begin{abstract}
Companion XII defines the relaxation scale geometrically as the comoving
horizon wavenumber at freeze-out,
\(k_{\mathrm{rel}}=a_{\mathrm{fr}}H(a_{\mathrm{fr}})\). This note provides
the numerical closure of that freeze-out relation and of the same scale
subsequently used in Companions L, XV, XLIV and the RDCC--Euclid Interface.
Evaluating the product on the Phase-1 background of Companion XI, with
\(\Omega_m\simeq0.31\), \(\Omega_r h^2\simeq4.15\times10^{-5}\) and
\(h\simeq0.67\), one finds
\(k_{\mathrm{rel}}\simeq0.037\,h\,\mathrm{Mpc}^{-1}\) for the benchmark
\(a_{\mathrm{fr}}=10^{-4}\), and
\(k_{\mathrm{rel}}\simeq0.015\,h\,\mathrm{Mpc}^{-1}\) if freeze-out coincides
with matter--radiation equality. Both values lie within a factor of a few of
the historical working scale \(0.01\,h\,\mathrm{Mpc}^{-1}\). The residual
uncertainty is dominated by \(a_{\mathrm{fr}}\), not by \(H(a)\). The
relaxation correction \(\delta_{\mathrm{rel}}=\alpha_{\mathrm{IR}}f_{\mathrm{rel}}\)
shifts the result only at the percent level.
\end{abstract}

\tableofcontents
\newpage

\section{Introduction}

This note does not redefine freeze-out and does not introduce a new
dynamical mechanism. It answers one question: given the geometric
definition of Companion XII and the background of Companion XI, what
number is \(k_{\mathrm{rel}}\)?

The same \(k_{\mathrm{rel}}\) is the transition scale in the growth sector
(Companions VII and L), the gravitational slip (Companion XLIV), tensor
damping (Companion XV) and the Euclid interface. A numerical value is
therefore the closure of that chain, not an extra parameter.

\section{Definition}

Companion XII identifies
\begin{equation}
  k_{\mathrm{rel}}
  =
  a_{\mathrm{fr}} H(a_{\mathrm{fr}}),
  \label{eq:def}
\end{equation}
with \(a_0=1\). Equation~\eqref{eq:def} is a comoving wavenumber. The Hubble
parameter must therefore be converted with
\begin{equation}
  \frac{H_0}{c}
  =
  3.3356\times10^{-4}\, h\,\mathrm{Mpc}^{-1}.
\end{equation}

\section{Background from Companion XI}

The working expansion history of Companion XI is
\begin{equation}
  H^2(a)
  =
  H_0^2
  \Bigl[
    \Omega_r a^{-4}
    +
    \Omega_m a^{-3}
    +
    \Omega_\Lambda
    +
    \delta_{\mathrm{rel}}(a)
  \Bigr],
\end{equation}
where \(\Omega_m=\Omega_b+\Omega_{\mathrm{DM,eff}}\) and
\begin{equation}
  \delta_{\mathrm{rel}}(a)
  =
  \alpha_{\mathrm{IR}}\, f_{\mathrm{rel}}(a).
\end{equation}
With \(\alpha_{\mathrm{IR}}=0.0176\pm0.0021\), \(\delta_{\mathrm{rel}}\) is an
\(\mathcal{O}(2\%)\) correction at freeze-out and is omitted from the leading
table. Benchmark parameters:
\begin{equation}
  \Omega_m \simeq 0.31,
  \qquad
  h \simeq 0.67,
  \qquad
  \Omega_r h^2 \simeq 4.15\times10^{-5},
  \qquad
  \Omega_\Lambda \simeq 1-\Omega_m-\Omega_r.
\end{equation}
These imply
\begin{equation}
  \Omega_r \simeq 9.25\times10^{-5},
  \qquad
  a_{\mathrm{eq}}
  =
  \frac{\Omega_r}{\Omega_m}
  \simeq
  3.0\times10^{-4}.
\end{equation}
At \(a=10^{-4}\) radiation still dominates, but the matter term is already
comparable (\(\Omega_r a^{-4}\sim9\times10^{11}\),
\(\Omega_m a^{-3}\sim3\times10^{11}\)).

\section{Numerical table}

\begin{center}
\begin{tabular}{lcccc}
\toprule
\(a_{\mathrm{fr}}\) &
\(H(a_{\mathrm{fr}})/H_0\) &
\(a_{\mathrm{fr}}H/H_0\) &
\(k_{\mathrm{rel}}\,[h\,\mathrm{Mpc}^{-1}]\) &
comment \\
\midrule
\(5\times10^{-5}\) &
\(4.16\times10^{6}\) &
\(208\) &
\(0.069\) &
deep radiation era \\
\(1\times10^{-4}\) &
\(1.11\times10^{6}\) &
\(111\) &
\(0.037\) &
XII benchmark \\
\(2\times10^{-4}\) &
\(3.11\times10^{5}\) &
\(62\) &
\(0.021\) &
approaching equality \\
\(a_{\mathrm{eq}}\simeq3.0\times10^{-4}\) &
\(1.53\times10^{5}\) &
\(46\) &
\(0.015\) &
matter--radiation equality \\
\bottomrule
\end{tabular}
\end{center}

Including \(\delta_{\mathrm{rel}}\sim\alpha_{\mathrm{IR}}\) multiplies every
entry by \(1+\mathcal{O}(10^{-2})\) and does not change the table at the
quoted precision.

A plot of \(k_{\mathrm{rel}}(a_{\mathrm{fr}})\) against \(a_{\mathrm{fr}}\),
with markers at the four rows above and a horizontal line at the historical
value \(0.01\,h\,\mathrm{Mpc}^{-1}\), is deferred to a later figure; the
table already contains the same information.

\section{Comparison with the historical scale}

Earlier RDCC documents used
\begin{equation}
  k_{\mathrm{rel}}^{\mathrm{(hist)}}
  \sim
  0.01\,h\,\mathrm{Mpc}^{-1}
\end{equation}
as a working value in growth, RSD and Euclid forecasts. The geometric
evaluation yields
\begin{equation}
  k_{\mathrm{rel}}
  \simeq
  0.037\,h\,\mathrm{Mpc}^{-1}
  \qquad
  (a_{\mathrm{fr}}=10^{-4})
\end{equation}
and
\begin{equation}
  k_{\mathrm{rel}}
  \simeq
  0.015\,h\,\mathrm{Mpc}^{-1}
  \qquad
  (a_{\mathrm{fr}}=a_{\mathrm{eq}}).
\end{equation}
The mismatch is a factor \(\sim2\)--\(4\), not an order of magnitude. The
historical number is essentially \(k_{\mathrm{eq}}\). The theory therefore
does not fall out of the observational window. The dominant uncertainty is
whether freeze-out sits at \(10^{-4}\) or nearer equality, that is, whether
\(\beta(a_{\mathrm{fr}})=1\) is realized before or at \(a_{\mathrm{eq}}\).

\section{Consequences}

\begin{itemize}[leftmargin=0.8cm]
  \item \textbf{Growth and Companion L.}
        The turnover in \(D\), \(f\) and \(\gamma\) moves from
        \(0.01\) to \(\sim0.02\)--\(0.04\,h\,\mathrm{Mpc}^{-1}\) unless
        \(a_{\mathrm{fr}}\) is revised toward \(a_{\mathrm{eq}}\).
  \item \textbf{Slip and Companion XLIV.}
        The scale where \(\eta=\Phi/\Psi\) departs from unity is the same
        number.
  \item \textbf{Tensors and Companion XV.}
        High-frequency SGWB suppression begins at this \(k_{\mathrm{rel}}\),
        not at an independent tensor scale. The bounce scale \(k_b\) remains
        distinct.
  \item \textbf{Euclid.}
        Forecast templates that inserted \(0.01\,h\,\mathrm{Mpc}^{-1}\) should
        treat that value as a prior center, not as the XII prediction, until
        \(a_{\mathrm{fr}}\) is derived from \(\beta(a)=1\).
\end{itemize}

\section{What this note does not decide}

Equation~\eqref{eq:def} is evaluated, not re-derived. The note does not
construct \(c_{\mathrm{eff}}\), does not fix \(\beta(a)\), and does not
resolve the two expressions for \(G_{\mathrm{eff}}\) in Companions I and XI.
Those tasks are separate. The immediate scientific content is the table:
the geometric scale lands in the same decade as the old working value, and
the remaining lever is \(a_{\mathrm{fr}}\).

\end{document}
